\documentclass[11pt]{article}
\usepackage[margin=1in]{geometry}
\usepackage{amsmath,amssymb,amsthm}
\usepackage[T1]{fontenc}

\newtheorem{theorem}{Theorem}
\newtheorem{lemma}{Lemma}
\newtheorem{corollary}{Corollary}
\newtheorem{definition}{Definition}

\title{V114: A One-Extra Shared-Gate Barrier for Signed-MUX Return Selection}
\author{NC0\_k-Avoid Laboratory}
\date{2026-08-21}

\begin{document}
\maketitle

\paragraph{Scope.}
This note formalizes the V114 route-level barrier.  It does not prove
NP-hardness of the V113 parameter $\Delta$, does not solve unrestricted
$NC^0_3$ range avoidance, and does not resolve P versus NP.

\begin{definition}[One-extra opposite-compatible return]
Fix a signed ternary MUX circuit, a root selector $r$, two distinct outputs
with selector $r$, and prescribed first branches whose selector-source phases
are opposite.  Remove the other selector-$r$ outputs from the return graph.
A certificate consists of two gate-simple return routes from the two first
branch destinations back to $r$ such that:
\begin{enumerate}
\item exactly one return output gate is used by both routes;
\item the two routes use opposite MUX branches at that shared gate; and
\item the target bit induced by the two implication cycles is the same at the
shared gate.
\end{enumerate}
The associated decision problem is denoted
\textsc{One-Extra-Opposite-Compatible-Return}.
\end{definition}

\begin{theorem}[V114 barrier theorem]
\textsc{One-Extra-Opposite-Compatible-Return} is NP-complete.  NP-hardness
holds for essential signed ternary MUX circuits with exact positive stretch
$m=n+1$.  Moreover, every instance produced by the reduction has, for the
same distinguished first pair, an explicit target-compatible return pair of
overlap zero.
\end{theorem}

\begin{proof}
Membership in NP is immediate: the two gate-simple routes have polynomial
length, and continuity, exact one-gate overlap, opposite branch bits, and the
induced shared target bit can all be checked in polynomial time.

For NP-hardness reduce from the directed two vertex-disjoint paths problem
with two ordered terminal pairs $(s_1,t_1)$ and $(s_2,t_2)$, which is
NP-complete by Fortune, Hopcroft, and Wyllie~\cite{FHW1980}.

Let $D=(V,A)$ be the directed input graph.  Create a fresh logical waypoint
$w$.  For every $v\in V$, create a variable $x_v$.  Unless $v=t_2$, also
create a private choice variable $c_v$.  Add a mandatory vertex-resource MUX
$Q_v$ whose selector is $x_v$ and whose branch zero goes to $c_v$; its other
branch goes to a dead trap.  For $v=t_2$, branch zero instead reaches the root
$r$.  For every arc $(u,v)\in A$ with $u\ne t_2$, add an arc MUX selected by
$c_u$ whose branch zero reaches $x_v$ and whose other branch reaches the dead
trap.  Add one unique connector from $c_{t_1}$ to $x_w$.

The waypoint gate $W$ is selected by $x_w$.  Its branch zero reaches
$x_{s_2}$ and its branch one reaches a private suffix leading to $r$.  Add a
root first pair whose prescribed branch-zero destinations are fresh variables
$a_0,a_1$, with selector-source phases respectively zero and one.  From $a_0$
provide a private graph-entry gate to $x_{s_1}$; from $a_1$ provide a private
prefix to $x_w$.  Finally give $a_0$ a second private route to $r$, disjoint
from the $a_1$ scaffold.  All unused branches enter a trap component that has
no route back to $r$.

If the raw construction has fewer than $n+1$ outputs, add essential ternary
MUX outputs entirely inside the dead component.  If it has more than $n+1$
outputs, add isolated unused variables.  Thus $m=n+1$ exactly, and the padding
cannot change return-route existence.

Suppose $D$ has vertex-disjoint directed paths $P_1:s_1\leadsto t_1$ and
$P_2:s_2\leadsto t_2$.  The first return route follows the encoded gates for
$P_1$, the unique $t_1\to w$ connector, branch zero of $W$, the encoded gates
for $P_2$, and finally $Q_{t_2}$ to $r$.  The second route follows its private
prefix to $W$, branch one of $W$, and its private suffix to $r$.  The routes
share exactly $W$.  The first gates after the two branches of $W$ both have
source phase zero, and the waypoint data polarities are zero, so both
traversals require the same target bit at $W$.  Hence the certificate is
accepted.

Conversely, consider any accepting certificate.  The private second source
has only one root-reaching prefix and therefore every second return uses $W$.
Every ordinary vertex-resource gate and every arc gate has only branch zero on
any root-reaching route.  Such a gate therefore cannot be the unique shared
gate of a certificate that requires opposite branch bits.  The shared gate is
necessarily $W$.

A return reaching $W$ from the first source without taking the private bypass
must alternate mandatory vertex-resource and encoded arc gates beginning at
$s_1$.  The only encoded entry to $w$ is the connector from $t_1$, so this
prefix determines a directed path $P_1:s_1\leadsto t_1$.  Whichever return
uses branch zero of $W$ then alternates encoded gates from $s_2$ until
$Q_{t_2}$ reaches $r$, determining a path $P_2:s_2\leadsto t_2$.  If both
graph portions occur on one return, gate simplicity prevents reuse of any
mandatory $Q_v$; if they occur on different returns, exact one-gate overlap
prevents both routes from using the same $Q_v$.  Thus $P_1$ and $P_2$ are
vertex-disjoint.  This proves equivalence.

The private first-source bypass and the private second-source scaffold have
disjoint return-gate sets, hence give an explicit compatible overlap-zero pair.
Therefore the hard waypoint is an extra non-dominator shared gate.  The
reduction is polynomial, completing NP-hardness and the theorem.
\end{proof}

\begin{corollary}[Generic guessed-waypoint barrier]
Unless $\mathrm{P}=\mathrm{NP}$, there is no unrestricted polynomial oracle
that, after choosing one arbitrary extra shared MUX gate and an opposite-branch
state, always decides the corresponding two-route completion problem.  This
corollary does not rule out an algorithm that exploits structural information
specific to rejection of the entire V113 minimum-overlap face.
\end{corollary}

\paragraph{Important non-implication.}
The reduction contains a compatible overlap-zero pair in every instance.
Consequently it does not establish NP-hardness for deciding whether
$\Delta=1$, and it does not establish parameterized hardness in $\Delta$.

\begin{thebibliography}{9}
\bibitem{FHW1980}
S.~Fortune, J.~Hopcroft, and J.~Wyllie,
\emph{The directed subgraph homeomorphism problem},
Theoretical Computer Science 10 (1980), 111--121.
\end{thebibliography}

\end{document}
